%!TEX TS-program = xelatex
%!TEX encoding = UTF-8 Unicode
\documentclass[manuscript,screen,review]{acmart}

\setcopyright{none}
\copyrightyear{2026}
\acmYear{2026}
\acmJournal{TOSEM}
\settopmatter{printacmref=true, printccs=true, printfolios=true}
\setcitestyle{numbers,sort&compress,square}

\usepackage{amssymb,amsthm}
\usepackage{mathtools}
\usepackage{calc}
\usepackage{enumitem}
\usepackage{longtable}
\usepackage{array}
\usepackage{listings}
\usepackage{pifont}

\setlength{\emergencystretch}{3em}
\lstset{
  breaklines=true,
  breakatwhitespace=true,
  basicstyle=\ttfamily\small,
  frame=none,
  columns=fullflexible
}

\title[Supplementary Material]{Supplementary Material for ``A Semantic Mutation Metric for Metamorphic-Relation Adequacy in Scientific Computing Programs''}

\author{Meng Li}
\orcid{0000-0002-1074-1502}
\email{mlemon@usc.edu.cn}
\affiliation{%
  \institution{School of Computing, University of South China}
  \city{Hengyang}
  \country{China}
}
\affiliation{%
  \institution{Hunan Engineering Research Center of Software Evaluation and Testing for Intellectual Equipment}
  \city{Hengyang}
  \country{China}
}
\affiliation{%
  \institution{CNNC Key Laboratory on High Trusted Computing}
  \city{Hengyang}
  \country{China}
}

\author{Xiaohua Yang}
\email{xiaohua1963@foxmail.com}
\affiliation{%
  \institution{School of Computing, University of South China}
  \city{Hengyang}
  \country{China}
}

\author{Jie Liu}
\email{jieliu@usc.edu.cn}
\affiliation{%
  \institution{School of Computing, University of South China}
  \city{Hengyang}
  \country{China}
}

\author{Shiyu Yan}
\orcid{0000-0001-7626-5185}
\email{yanshiyu@usc.edu.cn}
\affiliation{%
  \institution{School of Computing, University of South China}
  \city{Hengyang}
  \country{China}
}

\begin{CCSXML}
<ccs2012>
 <concept>
  <concept_id>10011007.10011074.10011099.10011102</concept_id>
  <concept_desc>Software and its engineering~Software testing and debugging</concept_desc>
  <concept_significance>500</concept_significance>
 </concept>
 <concept>
  <concept_id>10011007.10011074.10011099</concept_id>
  <concept_desc>Software and its engineering~Software verification and validation</concept_desc>
  <concept_significance>300</concept_significance>
 </concept>
</ccs2012>
\end{CCSXML}

\ccsdesc[500]{Software and its engineering~Software testing and debugging}
\ccsdesc[300]{Software and its engineering~Software verification and validation}

\keywords{metamorphic testing, mutation testing, semantic mutation operators, metamorphic relation adequacy, LLM-generated mutants, scientific computing}

\begin{document}

\maketitle

\appendix

\section{A. Notation and Operator Catalogue}\label{a.-notation-and-operator-catalogue}

\subsection{A.1 Complete notation table}\label{a.1-complete-notation-table}

\paragraph{Experimental objects.}
\begin{itemize}
\item PUT set $I = \{\text{A1,A2,A3, B1,B2,B3, C1,C2,C3, D1,D2,D3}\}$, $|I| = 12$.
\item PUT: $S_i$ ($i \in I$).
\item Class mapping: $\mathrm{cls} : I \to \{A, B, C, D\}$ (open, extensible).
\item Valid input distribution: $D_S$, sampling $X_{K_{\mathrm{eq}}} \sim D_S$.
\end{itemize}

\paragraph{Metamorphic testing side.}
\begin{itemize}
\item Meta-Pattern (MP): $\mathrm{MP} = \{\mathrm{MP}_1, \ldots, \mathrm{MP}_5\}$, $k \in \{1,\ldots,5\}$ (open, extensible). MP\_1 Conservation, MP\_2 Monotonicity, MP\_3 Convergence, MP\_4 Trajectory, MP\_5 Partial-order.
\item Metamorphic Relation: $\mathrm{MR}_{i,k}$ (from the shared infrastructure); $\mathit{mr} = (r, R) \in \mathrm{MR}_{i,k}$.
\item Automated Verification Pipeline (AVP): $\mathrm{AVP} : \text{Programs} \times \mathrm{MR}_{\text{universe}} \times \mathbb{R}^+ \to \{\text{pass}, \text{fail}\}$; $\mathrm{AVP}(s, \mathit{mr}, \varepsilon_{\mathrm{AVP}}^k)$ invokes the verification method per MP\_$k$.
\end{itemize}

\paragraph{Mutation testing side.}
\begin{itemize}
\item Mutation operator family: $\mathrm{MUT} = \{\mathrm{mut}_1, \ldots, \mathrm{mut}_5\}$, $j \in \{1,\ldots,5\}$ (open, extensible); $\mathrm{mut}_1 = \mathrm{mut}_C$, $\mathrm{mut}_2 = \mathrm{mut}_M$, $\mathrm{mut}_3 = \mathrm{mut}_G$, $\mathrm{mut}_4 = \mathrm{mut}_T$, $\mathrm{mut}_5 = \mathrm{mut}_F$.
\item Signature: $\mathrm{mut}_j : \text{Programs} \to 2^{\text{Programs}}$.
\item Mutant: $s' \in \mathrm{mut}_j(S_i)$.
\item Alignment: $\mathrm{align}(j) = j$.
\end{itemize}

\paragraph{Three-state decomposition (fixed).}
\[
\begin{aligned}
\mathrm{mut}_j(S_i)
  &= \mathrm{equiv}_{i,k,j} \cup \mathrm{killed}_{i,k,j}
     \cup \mathrm{survive}_{i,k,j},\\
&\hspace{1.5em}\text{with the three sets disjoint, mutually exclusive, and exhaustive.}
\end{aligned}
\]
\begin{itemize}
\item \emph{equiv determination} (dual conditions):
  (E1) AVP-coherent: $\forall \mathit{mr} \in \mathrm{MR}_{i,k}: \mathrm{AVP}(S_i, \mathit{mr}) = \mathrm{AVP}(s', \mathit{mr})$;
  (E2) Output-equiv: $\forall x \in X_{K_{\mathrm{eq}}} \sim D_S: \|S_i(x) - s'(x)\| \leq \varepsilon_{\mathrm{eq}}$.
\item \emph{killed determination}: $\mathrm{killed}(s', \mathrm{MR}_{i,k}) \Leftrightarrow \exists \mathit{mr} \in \mathrm{MR}_{i,k}: \mathrm{AVP}(S_i, \mathit{mr}) = \text{pass} \wedge \mathrm{AVP}(s', \mathit{mr}) = \text{fail}$.
\end{itemize}

\paragraph{Metric (fixed classical structure).}
\begin{align*}
\mathrm{SMS}_{i,k,j} &:= \frac{|\mathrm{killed}_{i,k,j}|}{|\mathrm{mut}_j(S_i)| - |\mathrm{equiv}_{i,k,j}|} \\
                    &\phantom{:}= \frac{|\mathrm{killed}_{i,k,j}|}{|\mathrm{killed}_{i,k,j}| + |\mathrm{survive}_{i,k,j}|} \in [0, 1].
\end{align*}

\paragraph{LRCA engineering attribution layer (descriptive, not in SMS).}
\begin{itemize}
\item Likely root-cause inventory: $C = \{C_1, \ldots, C_5\}$ (open, extensible). C1 True semantic failure; C2 Tolerance perturbation; C3 OOD; C4 Statistical-assumption violation; C5 Mutator artefact.
\item LRCA output: each $s' \in \mathrm{killed}$ annotated with $\texttt{root\_cause}(s') \in C$.
\item Descriptive quantities: $\texttt{C1\_share}_{i,k,j}$, $\texttt{suspect\_share}_{i,k,j}$.
\end{itemize}

\paragraph{Slicing and cross-class.}
\begin{itemize}
\item Slicing: aligned $j = k$; cross $j \neq k$.
\item Cross-class: $\Delta\mathrm{SMS}_c$, $c \in \{A, B, C, D\}$; $\mathrm{CV}(\Delta\mathrm{SMS}) := \mathrm{std}(\Delta\mathrm{SMS}_c) / |\mathrm{mean}(\Delta\mathrm{SMS}_c)|$.
\end{itemize}

\paragraph{Tolerance and sampling.}
\begin{itemize}
\item Tolerance: $\varepsilon_{\mathrm{AVP}}^k$ (MP-dependent), $\varepsilon_{\mathrm{eq}}$ (equivalence output tolerance).
\item Sampling: $K_{\mathrm{eq}} = 1000$ (E2 sample size), $N = 20$ (statistical replicates).
\end{itemize}

The notation skeleton (11 core concepts + three-state decomposition +
SMS formula + AVP interface) is \textbf{fixed}; the \emph{content} of
MUT, MP, C, and cls is open to extension.

\subsection{A.2 LRCA decision tree and three-layer diagnostic protocol}\label{a.2-lrca-decision-tree-and-three-layer-diagnostic-protocol}

\textbf{Five likely root causes} (open):

{%
\begin{longtable}[]{@{}ll@{}}
\caption{LRCA decision tree and three-layer diagnostic protocol.}\label{tab:p2-15}\\
\toprule\noalign{}
Code & Meaning \\
\midrule\noalign{}
\endfirsthead
\endhead
\bottomrule\noalign{}
\endlastfoot
C1 & True semantic failure (what SMS seeks) \\
C2 & Numerical tolerance perturbation \\
C3 & Out-of-distribution (OOD) \\
C4 & Statistical-assumption violation \\
C5 & Mutator artefact \\
\end{longtable}
}

\textbf{Three-layer diagnosis with L0 artefact pre-scan:}

\begin{itemize}
\item \textbf{L0 artefact pre-scan.} Double-blind review (§3.9.4 dual-LLM cross-source + 20\% manual sampling).
\item \textbf{L1 tolerance robustness.} $N = 20$ replicates; fail ratio $\geq 0.80$ considered stable.
\item \textbf{L2 OOD triage.} For C/D classes, distinguish valid $D_S^{\text{valid}}$ from OOD.
\item \textbf{L3 statistical baseline.} For Wilcoxon/DTW, IID/stationarity pre-check on PUT samples.
\end{itemize}

\textbf{Decision tree:}

For each $s' \in \mathrm{killed}_{i,k,j}$:
\begin{itemize}
\item \textbf{L1.} fail ratio $< 0.80 \Rightarrow$ C2; otherwise $\Rightarrow$ L2.
\item \textbf{L2} (C/D classes). fail only in OOD $\Rightarrow$ C3; otherwise $\Rightarrow$ L3.
\item \textbf{L3} (B/D classes + Wilcoxon/DTW). assumption violation $\Rightarrow$ C4; otherwise $\Rightarrow$ artefact recheck.
\item \textbf{Artefact recheck.} artefact evidence $\Rightarrow$ C5; otherwise $\Rightarrow$ C1.
\end{itemize}

\textbf{Multi-label priority:} C5 \textgreater{} C4 \textgreater{} C3
\textgreater{} C2 \textgreater{} C1.

\textbf{LRCA threshold calibration (9-grid scan;
\texttt{lrca\_calibration.json}).} Best ood\_band=0.02 with any
tolerance\_multiplier (3.0 / 10.0 / 30.0) gives mean\_C1\_share 0.200 /
H4 12/60 (20.0\%); ood\_band=0.05 (default) or 0.10 gives 0.164 / 10/60
(16.7\%). \texttt{tolerance\_multiplier} has zero impact (L1 tolerance
rarely triggered); \texttt{ood\_band} is the only discriminator.
Calibration ceiling 12/60 = 20\% remains far below 80\%, H4
unattainable on this dataset, not a calibration issue (inherent SNR of
LLM-mutant pools).

\subsection{A.3 AVP protocol specification and equivalence-judgement trade-off}\label{a.3-avp-protocol-specification-and-equivalence-judgement-trade-off}

\[
\mathrm{AVP}: \mathrm{Programs} \times \mathrm{MR}_{\mathrm{universe}} \times \mathbb{R}^{+} \to \{\mathrm{pass}, \mathrm{fail}\}.
\]
Per-MP verification: $\mathrm{MP}_1$ tolerance equality $|\mathrm{LHS} -
\mathrm{RHS}| \leq \varepsilon$; $\mathrm{MP}_2$ / $\mathrm{MP}_5$ Wilcoxon signed-rank $\alpha=0.05$; $\mathrm{MP}_3$
convergence-order + asymptotic residual ratio; MP\_4 DTW distance
threshold $\varepsilon_{\mathrm{DTW}}$. AVP version = upstream infrastructure commit hash; our package embeds
AVP source.

\textbf{Three-candidate equivalence trade-off.}

{%
\begin{longtable}[]{@{}
  >{\raggedright\arraybackslash}p{(\linewidth - 6\tabcolsep) * \real{0.2500}}
  >{\raggedright\arraybackslash}p{(\linewidth - 6\tabcolsep) * \real{0.2500}}
  >{\raggedright\arraybackslash}p{(\linewidth - 6\tabcolsep) * \real{0.2500}}
  >{\raggedright\arraybackslash}p{(\linewidth - 6\tabcolsep) * \real{0.2500}}@{}}
\caption{AVP protocol specification and equivalence-judgement trade-off.}\label{tab:p2-16}\\
\toprule\noalign{}
\begin{minipage}[b]{\linewidth}\raggedright
Determination
\end{minipage} & \begin{minipage}[b]{\linewidth}\raggedright
False-positive
\end{minipage} & \begin{minipage}[b]{\linewidth}\raggedright
False-negative
\end{minipage} & \begin{minipage}[b]{\linewidth}\raggedright
SMS bias
\end{minipage} \\
\midrule\noalign{}
\endfirsthead
\endhead
\bottomrule\noalign{}
\endlastfoot
E1 alone & numerical coincidence + AVP-coherent & E1 false $\to$ non-equiv &
Low (smaller denominator) \\
E2 alone & output match, MR differs (rare) & K\_eq miss, full-space
equal & High \\
\textbf{E1 $\wedge$ E2} & both mis-pass simultaneously & E1 $\vee$ E2 false $\to$ non-equiv
& Slightly high (conservative) \\
\end{longtable}
}

E1 alone is fooled by insufficient AVP coverage; E2 alone by numerical
coincidence on $K_{\mathrm{eq}}$. E1 $\wedge$ E2 is the conservative complete
instantiation; under $L_{\mathrm{equiv}}$ it reduces almost-everywhere to
classical bitwise equivalence (Lemma 9.1). Counter-examples: E2 passes
but E1 does not (rare; mutant coincides at K\_eq but deviates at MR
trigger points $\to$ non-equivalent); E1 passes but E2 does not (common;
consistent within MR but numerical drift \textgreater{} $\varepsilon_{\mathrm{eq}} \to$
non-equivalent).

% [thematic break removed]

\section{B. Experimental Subjects and Mutation-Operator Specialisations}\label{b.-experimental-subjects-and-mutation-operator-specialisations}

\subsection{B.1 PUT selection rationale}\label{b.1-put-selection-rationale-detailed-coverage-argument}

\textbf{(a) Library-stack coverage.} numpy 2.4.4 (A1-B3 linear algebra /
arrays), scipy 1.17.1 (A1 integrate, A2 linalg, B2 stats), scikit-learn
1.8.0 (C1-D3 surrogate / ML), Python scientific-computing's de facto
foundation stack (PyPI top-50, 2026-04). Uncovered: GPU / distributed
(JAX, CuPy, dask), domain-specific (BioPython, Astropy, RDKit). Left to
future work (R5).

\textbf{(b) Mathematical-structure coverage.} ODE/PDE (A1, A3), direct
linear algebra (A2), Bayesian analytic + MCMC + Monte Carlo (B1-B3),
kernel + orthogonal polynomial + NN surrogate (C1-C3), convex
optimisation + backprop + max-likelihood (D1-D3), covers 8 of 12
chapters in Numerical Recipes \citep{press2007numerical}. Uncovered: advanced
PDE solvers (FEM, FV, spectral); FFT; interior-point/trust-region
optimisation; symbolic/CAS. Left to future work.

\textbf{(c) Comparison with benchmarks.} DeepCrime (Humbatova et
al.~2021): semantic class-D topical overlap on ML kernels. Defects4J (Just et
al.~2014): no overlap, only mutation-as-fault-proxy reference. mutmut /
cosmic-ray demos: general Python; the semantic operators extend the scientific-computing focus.

\textbf{(d) Scale vs representativeness.} Each PUT 50-400 LOC; signature
standardised to \texttt{program(x:\ float)\ $\to$\ float}. Substantively
smaller than industrial code (typically 1-10 KLOC).

\textbf{Limitation.} The \texttt{float\ $\to$\ float} signature is a
substantive constraint, not engineering trade-off. Industrial inputs are
typically high-dimensional (CFD grids, MD particles, FEM matrices);
scalarised PUTs upper-bound mutant semantic complexity and may
systematically under-estimate SMS and cross-class differences on
industrial PUTs (declared §6 R5; future work validates on industrial-grade PUTs).

\subsection{B.1.5 Systematic vs incidental}\label{b.1.5-systematic-vs-incidental}

Syntactic tools may \emph{occasionally} hit §3.2 (a)
(cross-function-boundary) or (c) (algorithm class), but occasionality
undermines two engineering functions of semantic mutation. \textbf{(a)
Deepening source-code understanding.} Designing OS
\texttt{np.linalg.det(M)\ $\to$\ np.sum(np.diag(M))} requires knowing the
two are equivalent on diagonal matrices but not on general matrices
(\texttt{det\ =\ $\prod$\ eigvals} vs
\texttt{sum(diag)\ =\ trace\ =\ $\sum$\ eigvals}); syntactic tool AST
traversal does not require this understanding even when similar
replacements appear. \textbf{(b) Revealing deep faults.}
Domain-semantic errors (physical-constant, unit conversion, boundary
conditions, hyperparameter semantics, numerical-method order) are not
AST-local; syntactic mutator design goals (operator typos, off-by-one,
negation flips) have hit probability for domain errors much smaller than
systematic semantic-mutator triggering, and lack repeatability.
\textbf{Conclusion.} Syntactic tools occasionally producing mutants
satisfying (a)(b)(c) is stochastic byproduct, not repeatable, not
carrying engineering value. Systematic semantic mutation requires
(a)(b)(c) to be design intent.

\subsection{B.2 Per-class operator specialisations}\label{b.2-per-class-operator-specialisations}

{%
\begin{longtable}[]{@{}
  >{\raggedright\arraybackslash}p{(\linewidth - 2\tabcolsep) * \real{0.5000}}
  >{\raggedright\arraybackslash}p{(\linewidth - 2\tabcolsep) * \real{0.5000}}@{}}
\caption{Per-class operator specialisations.}\label{tab:p2-17}\\
\toprule\noalign{}
\begin{minipage}[b]{\linewidth}\raggedright
Operator
\end{minipage} & \begin{minipage}[b]{\linewidth}\raggedright
Representative specialisations across PUTs
\end{minipage} \\
\midrule\noalign{}
\endfirsthead
\endhead
\bottomrule\noalign{}
\endlastfoot
\textbf{mut\_C} Conservation-breaking & A1 Lorenz: add $\varepsilon_{\mathrm{drift}}$ to RHS
(slow Hamiltonian drift); A2 LU: decomposition omits k+1-th multiplier
(breaks det conservation); B1 Beta-Bin: posterior omits normalisation;
C1 GPR: covariance omits positive-definite diagonal term; D1 MLP:
backprop omits one gradient term. \\
\textbf{mut\_M} Monotonicity-breaking & A3 FDM: $\Delta$t coefficient
occasionally negative; B2 MCMC: acceptance min(1, r) $\to$ min(0.95, r); C2
PCE: high-order coefficient sort inserts inversion; D2 SVM:
decision-function sign flips near boundary. \\
\textbf{mut\_G} Convergence-breaking & A1 Lorenz: RK4 $\to$ 1.5-order
hybrid; A3 FDM: 2nd-order difference $\to$ 1st-order; B3 MC: doubling sample
size does not 1/N-reduce variance; C3 NN-Surr: training-epoch
truncation. \\
\textbf{mut\_T} Trajectory-distorting & A1 Lorenz: state-vector y/z
swap; B2 MCMC: insert independent-sampling segment; C3 NN-Surr:
training-target slow phase shift; D1 MLP: hidden-layer periodic-mask
activation. \\
\textbf{mut\_F} Fidelity-order-breaking & A2 LU: partial pivoting
degrades to no pivoting; C1 GPR: length-scale switches to coarse prior;
C2 PCE: high-order term randomly retains low-order; D3 LR:
regularisation occasionally large. \\
\end{longtable}
}

\textbf{Per-class HP / OS / TF substitution rules} (illustrative).
\textbf{HP} on class a: tolerance / max\_iter; on class c: GPR
\texttt{noise\_level} / \texttt{length\_scale}; on class d: MLP
\texttt{hidden\_dim} / dropout. \textbf{OS} on class a: numerical-linalg
API swaps (\texttt{det} $\leftrightarrow$ \texttt{sum(diag)}); on class b:
probability-distribution sampling swaps; on class c: surrogate-class
swaps (GPR $\leftrightarrow$ RBF $\leftrightarrow$ NN). \textbf{TF} on class a: integration-order
changes (RK4 $\to$ Euler); on class b: MC estimator changes.

\textbf{Operator-level cross-table with cosmic-ray defaults}. The 12
cosmic-ray default operators are AST-local:
\texttt{NumberReplacer} (Num/Constant, CE partial, $\bigtriangleup$ literal values
only); \texttt{ReplaceArithmeticOperator} (BinOp);
\texttt{ReplaceComparisonOperator} (Compare);
\texttt{ReplaceLogicalOperator} (BoolOp); \texttt{ReplaceUnaryOperator}
(UnaryOp); \texttt{ReplaceTrueFalse} (NameConstant, CE keyword, $\bigtriangleup$
Bool literal); \texttt{BreakContinueReplacer} (Break/Continue);
\texttt{RemoveDecorator}; \texttt{RemoveExceptHandler};
\texttt{ZeroIterationForLoop} (For); \texttt{ReplaceIfBlock} (If);
\texttt{MutateSubscript}. None correspond to OS / HP / TF / SI:

{%
\begin{longtable}[]{@{}
  >{\raggedright\arraybackslash}p{(\linewidth - 4\tabcolsep) * \real{0.3333}}
  >{\raggedright\arraybackslash}p{(\linewidth - 4\tabcolsep) * \real{0.3333}}
  >{\raggedright\arraybackslash}p{(\linewidth - 4\tabcolsep) * \real{0.3333}}@{}}
\caption{Per-class operator specialisations.}\label{tab:p2-18}\\
\toprule\noalign{}
\begin{minipage}[b]{\linewidth}\raggedright
Semantic class
\end{minipage} & \begin{minipage}[b]{\linewidth}\raggedright
Tool support
\end{minipage} & \begin{minipage}[b]{\linewidth}\raggedright
Coverage
\end{minipage} \\
\midrule\noalign{}
\endfirsthead
\endhead
\bottomrule\noalign{}
\endlastfoot
\textbf{OS} API replacement & None & $\bigtriangleup$ 88.33\% disjoint (incidental
low-complexity hits) \\
\textbf{HP} & None & $\times$ tool inexpressible (0/72) \\
\textbf{TF} numerical-method order & None & $\times$ tool inexpressible
(0/54) \\
\textbf{SI / CF} structural injection & None & $\times$ tool inexpressible
(0/33) \\
\end{longtable}
}

\textbf{Operator-level conclusion.} All 12 default classes remain
AST-local (BinOp / Compare / BoolOp / UnaryOp / NameConstant / Subscript
/ If / For / Break / Continue / decorator / except). Categorically, no
entry recognises sklearn / scipy hyperparameter semantics (HP),
numerical-method order (TF), or control-flow intent (SI / CF). For OS,
low-complexity sub-expressions can be incidentally hit by BinOp;
empirical 11.67\% (§3.5) is consistent with 88.33\% AST-disjointness
lower bound.

\subsection{B.3 cosmic-ray a1 single-PUT pilot}\label{b.3-cosmic-ray-a1-single-put-pilot-pre-12-put-empirical}

Before the 12-PUT generalisation (§3.5 main), a single-PUT lightweight
empirical used \texttt{scripts/run\_cosmic\_ray\_a1.sh} on a1 (file
934 B, AST-bound). Outcome: mutants\_generated \textasciitilde tens; $\geq$
90\% belong to BinOp / Compare / 12 default classes; small
killed-by-\texttt{test\_a1.py} ratio (PUT unit tests are
output-shape/type sanity checks; most cosmic-ray mutants not killed), an
empirical manifestation of the §3.2.6 thesis that syntactic-tool mutants
are not aligned with MR-violation detection. The 12-PUT generalisation
(§3.5) supersedes this pilot but corroborates the same conclusion at
scale.

\subsection{B.4 Expected LRCA risk profile}\label{b.4-expected-lrca-risk-profile}

\textbf{PUT-class $\times$ LRCA-layer risk weights.} A numeric: C2 dominant
($\star\star\star$ on L1 tolerance). B probabilistic: C4 dominant ($\star\star\star$ on L3). C
surrogate: C3 dominant ($\star\star\star$ on L2 OOD; $\star\star$ tolerance). D ML: C3 + C4
mixed ($\star\star\star$ L2; $\star\star$ L3).

\textbf{Operator-PUT root-cause hotspots (expected).} A: mut\_C/G/F
dominantly C2 (numerical tolerance), mut\_M/T C1 (true semantic). B:
dominantly C4 (statistical-assumption violations). C: mut\_C/M
dominantly C2/C1 on c1/c2, mut\_T/F dominantly C3 on c1/c2; c3 (NN-Surr)
row dominantly C5 (artefact). D: d1 row dominantly C5 (training-noise
artefact); d2/d3 dominantly C1 / C3 mixed.

\textbf{Expected suspect\_share thresholds.} A: 0.10-0.20 (acceptance $\leq$
0.25); B: 0.20-0.35 ($\leq$ 0.40); C: 0.20-0.30 ($\leq$ 0.35); D: 0.25-0.40 ($\leq$
0.45). 60-cell average: $\leq$ 0.20 (= H4; acceptance $\leq$ 0.25).

\subsection{B.5 Engineering-significance mapping}\label{b.5-engineering-significance-mapping}

j = k aligned diagonal is the H2 threshold-test slice; j $\neq$ k
off-diagonal is control; 6 vacant cells ($\circ$, no MR exercised) are not formally
adjudicated and are repurposed in §5.2 as descriptive evidence for
R\_sem / R\_kill decoupling. The v3 $\to$ v4 micro-change ($-$0.009)
attributes to LLM-source diversity under a fixed prompt; this is the
source-axis null contrast underpinning §4.3.

\subsection{B.6 Mutmut vs cosmic-ray default-operator overlap}\label{b.6-mutmut-vs-cosmic-ray-default-operator-overlap}

Reviewer-requested manual operator-class cross-reference between
cosmic-ray's 13 default operators (cosmic-ray 8.4.6,
\texttt{cosmic\_ray.operators} package) and mutmut's 14 default
operators (mutmut current \texttt{main},
\texttt{src/mutmut/mutation/mutators.py}). Both default-operator sets
are first-order AST-local replacements; neither implements a
cross-function-boundary, domain-knowledge-dependent, or
algorithmic-class-changing mutation (the three §3.2 necessary conditions
for semantic mutation).

{\scriptsize
\begin{longtable}[]{@{}
  >{\raggedright\arraybackslash}p{(\linewidth - 8\tabcolsep) * \real{0.1800}}
  >{\raggedright\arraybackslash}p{(\linewidth - 8\tabcolsep) * \real{0.2200}}
  >{\raggedright\arraybackslash}p{(\linewidth - 8\tabcolsep) * \real{0.2200}}
  >{\raggedright\arraybackslash}p{(\linewidth - 8\tabcolsep) * \real{0.1600}}
  >{\raggedright\arraybackslash}p{(\linewidth - 8\tabcolsep) * \real{0.2200}}@{}}
\caption{Mutmut vs cosmic-ray default-operator overlap.}\label{tab:p2-19}\\
\toprule\noalign{}
\begin{minipage}[b]{\linewidth}\raggedright
Operator class
\end{minipage} & \begin{minipage}[b]{\linewidth}\raggedright
cosmic-ray default
\end{minipage} & \begin{minipage}[b]{\linewidth}\raggedright
mutmut default
\end{minipage} & \begin{minipage}[b]{\linewidth}\raggedright
Reaches §3.2 (a/b/c)?
\end{minipage} & \begin{minipage}[b]{\linewidth}\raggedright
Semantic-class reachability
\end{minipage} \\
\midrule\noalign{}
\endfirsthead
\endhead
\bottomrule\noalign{}
\endlastfoot
Numeric literal $\pm$1 & \texttt{number\_replacer} &
\texttt{operator\_number} & (a) no, (b) no, (c) no & partial CE only
(e.g.~\texttt{\_RHO=28.0\ $\to$\ 27.5}) \\
Binary arithmetic swap (+, $-$, $\times$, $\div$) &
\texttt{binary\_operator\_replacement} & \texttt{operator\_swap\_op}
(binary subset) & (a) no, (b) no, (c) no & partial OS only
(incidental) \\
Comparison swap (\textless, $\leq$, \textgreater, $\geq$, ==, !=) &
\texttt{comparison\_operator\_replacement} & \texttt{operator\_swap\_op}
(compare subset) & (a) no, (b) no, (c) no & none \\
Boolean swap (and $\leftrightarrow$ or) & \texttt{boolean\_replacer} &
\texttt{operator\_swap\_op} (bool subset) & (a) no, (b) no, (c) no &
none \\
Unary remove / swap & \texttt{unary\_operator\_replacement} &
\texttt{operator\_remove\_unary\_ops}, \texttt{operator\_swap\_op} & (a)
no, (b) no, (c) no & none \\
Keyword swap (True/False, etc.) & \texttt{keyword\_replacer} &
\texttt{operator\_keywords}, \texttt{operator\_name} & (a) no, (b) no,
(c) no & none \\
break $\leftrightarrow$ continue / return & \texttt{break\_continue} &
\texttt{operator\_keywords} (subset) & (a) no, (b) no, (c) no & none \\
Exception class swap & \texttt{exception\_replacer} & (none) & (a) no,
(b) no, (c) no & none \\
Decorator removal & \texttt{remove\_decorator} & (none) & (a) no, (b)
no, (c) no & none \\
Variable-binding insert / replace & \texttt{variable\_inserter},
\texttt{variable\_replacer} & \texttt{operator\_arg\_removal},
\texttt{operator\_assignment} & (a) no, (b) no, (c) no & none \\
no\_op insertion & \texttt{no\_op} & (none) & (a) no, (b) no, (c) no &
none \\
Zero-iteration \texttt{for} loop & \texttt{zero\_iteration\_for\_loop} &
(none) & (a) no, (b) no, (c) no & none \\
String content tweak (XX prefix, case flip) & (none) &
\texttt{operator\_string} & (a) no, (b) no, (c) no & none \\
Lambda body $\to$ \texttt{None}/\texttt{0} & (none) &
\texttt{operator\_lambda} & (a) no, (b) no, (c) no & none \\
Dict kwarg name (XX prefix) & (none) &
\texttt{operator\_dict\_arguments} & (a) no, (b) no, (c) no & none \\
Sym./asym. string-method swap & (none) &
\texttt{operator\_symmetric\_string\_methods\_swap},
\texttt{operator\_unsymmetrical\_string\_methods\_swap} & (a) no, (b)
no, (c) no & none \\
Augmented $\to$ simple assignment & (none) &
\texttt{operator\_augmented\_assignment} & (a) no, (b) no, (c) no &
none \\
\texttt{match} case removal & (none) & \texttt{operator\_match} & (a)
no, (b) no, (c) no & none \\
\end{longtable}
}

\textbf{Aggregate.} Cosmic-ray $\cap$ mutmut: \textasciitilde9 of 13
cosmic-ray operators have a near-equivalent in mutmut (numeric, binary,
compare, boolean, unary, keyword, break-continue, variable, plus partial
overlap on string-prefix vs \texttt{XX}); 4 cosmic-ray operators are
unique (exception, decorator, no\_op, zero-iteration-for); 6 mutmut
operators are unique (string content, lambda, dict-kwarg, string-method
swap, aug-assignment, match-case). The union is approximately 21
distinct first-order AST-local operator classes. None of the 21 can:

\begin{itemize}

\item
  \begin{enumerate}[label=(\alph*)]
  
  \item
    Cross a function-call or module-import boundary (e.g.~replace
    \texttt{det(M)} with \texttt{sum(np.diag(M))} requires knowing the
    equivalence on diagonal matrices, outside any default-operator
    scope);
  \end{enumerate}
\item
  \begin{enumerate}[label=(\alph*)]
  \setcounter{enumi}{1}
  
  \item
    Depend on domain knowledge for legality (e.g.~a hyperparameter swap
    that knows \texttt{noise\_level=1e-4} is the load-bearing knob in a
    Gaussian Process Regression PUT requires PUT-class awareness);
  \end{enumerate}
\item
  \begin{enumerate}[label=(\alph*)]
  \setcounter{enumi}{2}
  
  \item
    Change the algorithmic class (e.g.~swap
    \texttt{polynomial\ $\to$\ spline\ basis} rewrites the surrogate's
    mathematical structure).
  \end{enumerate}
\end{itemize}

The mutmut / cosmic-ray null intersection on §3.5's HP/SI/TF classes
(zero overlap) therefore reflects a \textbf{structural property of
first-order AST-local mutation} rather than a tool-specific limitation,
supporting the §3.6 preventive-defence framing.

% [thematic break removed]

\section{C. Experimental Procedure Details}\label{c.-experimental-procedure-details}

\subsection{C.1 Cross-source protocol, full specification}\label{c.1-cross-source-v4-protocol-full-specification}

\textbf{Two-stage ablation.}

{%
\begin{longtable}[]{@{}
  >{\raggedright\arraybackslash}p{(\linewidth - 6\tabcolsep) * \real{0.2500}}
  >{\raggedright\arraybackslash}p{(\linewidth - 6\tabcolsep) * \real{0.2500}}
  >{\raggedright\arraybackslash}p{(\linewidth - 6\tabcolsep) * \real{0.2500}}
  >{\raggedright\arraybackslash}p{(\linewidth - 6\tabcolsep) * \real{0.2500}}@{}}
\caption{Cross-source protocol, full specification.}\label{tab:p2-20}\\
\toprule\noalign{}
\begin{minipage}[b]{\linewidth}\raggedright
Version
\end{minipage} & \begin{minipage}[b]{\linewidth}\raggedright
Mutant pool
\end{minipage} & \begin{minipage}[b]{\linewidth}\raggedright
c-class primary MP
\end{minipage} & \begin{minipage}[b]{\linewidth}\raggedright
Use
\end{minipage} \\
\midrule\noalign{}
\endfirsthead
\endhead
\bottomrule\noalign{}
\endlastfoot
v3 & Single-source (Claude Opus 4.6) & MP5 (pre-registered) &
H2 baseline \\
v4 & \textbf{Cross-source} (Claude Opus 4.6 + GPT-5.4 + DeepSeek chat) &
MP5 (held at pre-registered) & Isolate LLM-source diversity \\
\end{longtable}
}

\textbf{Protocol} (\texttt{scripts/cross\_source\_campaign.py}). (a) For
each (PUT, operator), Claude / GPT / DeepSeek run K = 3 trials with
identical prompt (§3.9.2), temperature 0.7; \texttt{source\_tag}
propagated to filenames. (b) \textbf{V1-V4 mechanical validation}
(\texttt{src/p2/mutators/validation.py}): V1 syntax
(\texttt{ast.parse}), V2 executability, V3 non-triviality
(\texttt{\textbar{}y\_mutant\ -\ y\_original\textbar{}\ \textgreater{}\ 1e-6}),
V4 signature consistency. (c) \textbf{DeepSeek model:}
\texttt{deepseek-chat} (113 tokens/call, 1.8 s), not
\texttt{deepseek-v4-pro} (340 tokens/call, 11 s), quality-equivalent
on a2\_OS1 dry-run. (d) \textbf{Pool capacity:} 37 ops $\times$ 3 sources $\times$ 3
trials = 333 attempts; V1-V4 pass rate 89.5\%; 298 confirmed mutants, of
which 292 enter the v4 pool after final-stage rejection of duplicates and
QA failures. Sources contribute nearly equally: Claude 101 / GPT 98 / DeepSeek 99
(Phase A engineering finding). (e) \textbf{Sampling}
(\texttt{scripts/build\_pools.py}, POOL\_VERSION = v4): per-PUT 30 max,
measured mean 24.3, range 10-30. c1 (GPR) has only 10 because c1\_HP1
/ c1\_CE1 have V1-V4 near-zero pass (WhiteKernel noise\_level 1e-4 $\to$
1e-1 perturbation has minimal output impact, triggers V3 non-trivial
failure, itself §5.2 evidence).

\textbf{Protocol-asymmetry (R13).} v4 does not invoke reviewer LLM (cost
/ speed priority; a follow-up study will rerun dual-blind on the v4
grid, \textasciitilde\$5-8 USD); v3 used the original Phase-1 dual-blind
(Claude gen + GPT-5.4 review + DeepSeek arbitration). A fraction of the
v3 $\to$ v4 source-axis shift may be slight v4 quality decline rather than
LLM-source diversity.

\subsection{C.1.1 Differential prompt protocol}\label{c.1.1-differential-prompt-protocol-future-work-commitment}

The v3 $\to$ v4 micro-change of $-$0.009 may reflect ``three LLMs under
identical prompt'' rather than ``true upper bound of LLM diversity''.
Separation requires \emph{one tailored differential prompt per LLM}
(skeleton: \texttt{scripts/run\_differential\_prompt.py}).

\textbf{Design (2$\times$3 factorial, within-(PUT, operator)).} Factor A:
prompt template, 3 levels: \textbf{V\_canonical} (control, identical to
v4), \textbf{V\_persona} (per-LLM identity: Claude ``numerical-analysis
editor''; GPT ``scientific-software refactorer''; DeepSeek ``library-API
substitution specialist''), \textbf{V\_cot} (Claude
\texttt{\textless{}thinking\textgreater{}} tags; GPT ``step by step'';
DeepSeek stepped-reasoning). Factor B, LLM source (Claude / GPT-5.4 /
DeepSeek chat). K = 3 per cell; total 37 $\times$ 3 $\times$ 3 $\times$ 3 = \textbf{999
trials}.

\textbf{Exit criteria.} If max(delta) $-$ min(delta) \textless{} 0.05: the
v3 $\to$ v4 $-$0.009 source-axis null is robust. If $\geq$ 0.05 with prompt variant
pushing delta past 0.474: H2 revised from ``not met'' to
``prompt-conditional met''. If $\geq$ 0.05 but all delta \textless{} 0.474:
prompt sensitivity exists but H2 unchanged.

\textbf{Resources.} API \textasciitilde\$18-30 USD; wall 3-4 h at
concurrency 4; V1-V4 only, no reviewer LLM.

\subsection{C.2 Manual pilot lineage, LLM-client configuration, and L0 prescreen}\label{c.2-manual-pilot-lineage-llm-client-configuration-and-l0-prescreen}

The original §3.9 protocol used 60\% multi-LLM consensus (Claude Opus
4.6 + GPT-5.4 + DeepSeek chat unanimous) plus 40\% manual injection by
scientific-computing researchers. Each (mut\_j, S\_i) pair carried one
prompt template containing PUT source, semantic intent of mut\_j,
prohibition against revealing specific MRs (avoids prompt leakage), and
output requirements (syntactically correct, executable, single-point
modification, diff \textless{} 10 lines). Reproducibility parameters:
temperature 0.3 in manual-pilot stage (v4 cross-source uses 0.7 per
§C.1); fixed seed; 5 candidates per prompt with 2-3 retained.

\textbf{Double-blind review (Scheme C).} Generator LLM-G (Claude Opus)
and reviewer LLM-R (GPT-5.4) are cross-source with strict role
separation. LLM-R sees only PUT original + mutant code, unaware of
mut\_j category / MR content / generator identity, outputting a triple
(syntactic, executable, semantic-fault-injected). Double-confirmed
mutants enter pool; inconsistencies enter manual arbitration ($\leq$ 10\%).
From the double-confirmed pool, 20\% is sampled for manual review;
manual-vs-LLM-R inconsistency \textgreater{} 10\% triggers full manual
downgrade. Same-source bias mitigation stratifies LLM-R by PUT class.
Prompt-injection mitigation disables RAG / web search on the API path.

\textbf{L0 pool prescreen.} Each mutant passes static syntax (linters /
type checkers), a unit self-test (simple inputs produce finite output),
and double-blind sign-off. Target: 30-50 candidates per cell $\to$ 10-15
retained after prescreen.

\subsection{C.4 LRCA three-layer execution}\label{c.4-lrca-three-layer-execution-full}

Decision tree, multi-label priority (C5 \textgreater{} C4 \textgreater{}
C3 \textgreater{} C2 \textgreater{} C1), and output (\texttt{C1\_share},
\texttt{suspect\_share}) are stated in §A.2; the full pseudocode is
reproduced there. LRCA does \textbf{not} modify SMS; killed set is not
filtered. §4.7 H4 numbers use best calibrated combination
(\texttt{ood\_band\ =\ 0.02,\ tolerance\_multiplier\ =\ 3.0,\ repeats\ =\ 20});
default-threshold results retained in \texttt{lrca\_60cell\_v3.json} as
control. Interface with §3.6 risk profile: L0 prescan $\leftrightarrow$ §C.2 dual-LLM
double-blind; L1 tolerance $\leftrightarrow$ N = 20 repetition subprocess; L2 OOD $\leftrightarrow$
input-distribution definition; L3 statistical baseline $\leftrightarrow$ pre-check
protocol.

\subsection{C.4.1 Pilot calibration}\label{c.4.1-pilot-calibration-37-operators-k-1020}

A 2026 Q3 operator-level pilot ran 12 PUTs $\times$ 37 operators (12
\texttt{is\_key=True} at K=20, 25 at K=10; 470 trials total). Pilot
precursor quantities: R\_sem (V1-V6 $\wedge$ operator\_match pass rate),
D\_impl (median pairwise AST + literal + identifier Jaccard distance
among confirmed mutants), R\_kill (proportion of confirmed mutants
killed by AVP under that PUT's primary MP). Aggregate (N=37): R\_sem
median 0.50 mean 0.468 (0/37 with R\_sem=0); D\_impl median 0.42 mean
0.392 (1/37 with D\_impl $\approx$ 0); R\_kill median 0.00 mean 0.189. All 37
operators produced $\geq$ 1 V1-V6 passing mutant (verifying H1).

\textbf{Key R\_sem / R\_kill decoupling pattern.} R\_sem $\geq$ 0.9 $\wedge$ R\_kill
= 0 combinations all on HP / TF / OS operators in c / d classes (e.g.,
c1\_CE1 noise 1e-4$\to$1e-1 with R\_sem 1.00, R\_kill 0.00; c3\_HP1
relu$\to$tanh, c3\_TF1 max\_iter 1000$\to$5, d1\_HP1 MLP $\alpha$, d3\_HP1 LR C).
R\_sem-high $\wedge$ R\_kill = 1 combinations all on CE / OS operators in a / b
classes (a2\_CE1 LU det, a2\_OS1 prod$\to$sum, b1\_OS1 $\alpha$/$\beta$ swap, d1\_TF1
label flip). This is empirical evidence for §5.2 R\_sem / R\_kill
decoupling at cell level.

% [thematic break removed]

\section{D. Statistical Analysis Details}\label{d.-statistical-analysis-details}

\subsection{D.1 SMS variants and construction lemmas}\label{d.1-sms-variants-and-construction-lemmas}

\textbf{SMS\_unfiltered (reviewer comparison metric):}

\[
\mathrm{SMS}^{\mathrm{unfiltered}}_{i,k,j} := \frac{|\mathrm{killed}_{i,k,j}|}{|\mathrm{mut}_j(S_i)| - |\mathrm{equiv}_{i,k,j}|}
\]

(Same form as the primary metric; here \texttt{killed} does not distinguish C1 from C2--C5.)

The appendix of the reproducibility package provides cell-by-cell
difference between SMS\_unfiltered and SMS; if relative difference
\textless{} 5\%, this confirms LRCA does not affect robustness of
primary conclusions.

\textbf{Primary-table construction.}

{%
\begin{longtable}[]{@{}
  >{\raggedright\arraybackslash}p{(\linewidth - 4\tabcolsep) * \real{0.3333}}
  >{\raggedright\arraybackslash}p{(\linewidth - 4\tabcolsep) * \real{0.3333}}
  >{\raggedright\arraybackslash}p{(\linewidth - 4\tabcolsep) * \real{0.3333}}@{}}
\caption{SMS variants and construction lemmas.}\label{tab:p2-21}\\
\toprule\noalign{}
\begin{minipage}[b]{\linewidth}\raggedright
RQ
\end{minipage} & \begin{minipage}[b]{\linewidth}\raggedright
Primary statistics
\end{minipage} & \begin{minipage}[b]{\linewidth}\raggedright
Reporting format
\end{minipage} \\
\midrule\noalign{}
\endfirsthead
\endhead
\bottomrule\noalign{}
\endlastfoot
RQ3 & inst\_rate, equiv\_rate, C1\_share, survive\_rate & 60-cell
heatmap + 4-class marginals \\
RQ4 & aligned-SMS vs cross-SMS; sparse $\circ$ vs dense \textbullet\textbullet equiv\_rate &
Cliff's delta + odds ratio + 95\% bootstrap CI \\
RQ5 & $\Delta$SMS\_c (c $\in$ \{A,B,C,D\}); CV($\Delta$SMS) & sign test (df=3) +
descriptive forest plot \\
RQ5 (desc.) & Spearman $\rho$ + Kendall $\tau$ for SMS vs PC & scatter + dual-metric
ranking comparison \\
\end{longtable}
}

\textbf{Multiple comparisons.} Cell-level claims across 60 cells use
Benjamini-Hochberg FDR, $\alpha_{\mathrm{FDR}} = 0.05$. N = 20 bootstrap intervals:
1,000-iteration 95\% CI (B = 10,000 for the headline H2 delta CI).

\subsection{D.2 Cliff's delta cutoff sensitivity and logit-transformation invariance}\label{d.2-cliffs-delta-cutoff-sensitivity-and-logit-transformation-invariance}

\textbf{H4 cutoff sensitivity.} Dense grid scan (cutoff
$\in$ \{0.05, 0.10, \ldots, 0.50\}, step 0.05) on v4 data
(\texttt{scripts/h5\_sensitivity.py},
\texttt{data/results/h5\_sensitivity\_v4.json}):

{%
\begin{longtable}[]{@{}lll@{}}
\caption{Cliff's delta cutoff sensitivity and logit-transformation invariance.}\label{tab:p2-22}\\
\toprule\noalign{}
cutoff & h5\_cells\_pass & h5\_pass\_ratio \\
\midrule\noalign{}
\endfirsthead
\endhead
\bottomrule\noalign{}
\endlastfoot
0.05 & 12 / 60 & 20.0\% \\
0.10 & 12 / 60 & 20.0\% \\
0.15 & 12 / 60 & 20.0\% \\
\textbf{0.20 (paper)} & \textbf{12 / 60} & \textbf{20.0\%} \\
0.25 & 12 / 60 & 20.0\% \\
0.30 & 12 / 60 & 20.0\% \\
0.35 & 12 / 60 & 20.0\% \\
0.40 & 12 / 60 & 20.0\% \\
0.45 & 13 / 60 & 21.7\% \\
0.50 & 13 / 60 & 21.7\% \\
\end{longtable}
}

\textbf{Conclusion.} H4 verdict (not met) is intrinsic data property,
independent of cutoff choice. v4 suspect\_share distribution is severely
bimodal: median 1.0, mean 0.79; 48 cross cells at suspect\_share $\approx$ 1, 12
aligned cells in low end; the {[}0.20, 0.80{]} interior is nearly empty.
Specific value 0.20 is \textbf{not load-bearing}.

\textbf{Logit-transformation invariance.} Cliff's delta is rank-based
(function of $U / (n_1 \cdot n_2)$), mathematically invariant under any strictly
monotone transformation (Romano 2006). Logit is strictly monotone on (0,
1), so $\delta_{\mathrm{logit}} \equiv \delta_{\mathrm{raw}}$ is a construction result. \textbf{Not
additional robustness evidence}; H2 robustness threats come from §4.2
zero-mass dominance, not metric-scale choice.

\subsection{D.3 Power analysis, plug-in and stipulated-alternative}\label{d.3-power-analysis-plug-in-and-stipulated-alternative}

\textbf{Plug-in bootstrap.} With-replacement sample from observed
aligned (n=12) and cross (n=48) v4 SMS pools; N\_sim = 5,000, seed = 42
(\texttt{scripts/compute\_rq2\_power.py}, \texttt{rq2\_power\_v4.json}):

{%
\begin{longtable}[]{@{}lll@{}}
\caption{Power analysis, plug-in and stipulated-alternative.}\label{tab:p2-23}\\
\toprule\noalign{}
Threshold & Interpretation & Power \\
\midrule\noalign{}
\endfirsthead
\endhead
\bottomrule\noalign{}
\endlastfoot
delta \textgreater{} 0.000 & Any-effect & \textbf{0.997} \\
delta \textgreater{} 0.147 & Small & 0.966 \\
delta \textgreater{} 0.330 & Medium & 0.759 \\
\textbf{delta \textgreater{} 0.474} & \textbf{Large (H2)} &
\textbf{0.423} \\
\end{longtable}
}

Sample-size sweep ($n_{\mathrm{aligned}} \in \{6, 12, \ldots, 60\}$, $n_{\mathrm{cross}} = 4 \times n_{\mathrm{aligned}}$):
power for delta \textgreater{} 0 reaches 0.974 at n=6, 0.996 at 12,
plateaus thereafter. RQ4 primary analysis has very low sample-size
requirement; the bottleneck is the effect-size boundary itself, not
sampling noise. Even at $\delta_{\mathrm{truth}} \approx 0.474$, sample has only
\textasciitilde42\% chance of detecting under plug-in; but this does
\textbf{not} mean ``insufficient power is the cause of H2 not being
met'', observed $\delta_{\mathrm{v3}} = 0.323 < 0.474$, observed
$\delta_{\mathrm{v4}} = 0.314 < 0.474$; increasing sample only narrows CI.

\textbf{Stipulated-alternative.} Implemented in
mixture-weight (\path{scripts/compute_rq2_power_stipulated.py},
N\_sim = 2,000): SMS has heavy ties at zero (45/60 = 0), raw shifting is
discontinuous (any $\varepsilon$ \textgreater{} 0 jumps delta from 0.314 to 0.74).
Mixture: aligned' = (probability w) sample from (observed\_aligned +
0.001) + (probability 1$-$w) sample from observed\_aligned, calibrate w so
E{[}delta{]} $\approx$ 0.474. Calibrated w = 0.094, realised E{[}delta{]} =
0.4746.

{%
\begin{longtable}[]{@{}lll@{}}
\caption{Power analysis, plug-in and stipulated-alternative.}\label{tab:p2-24}\\
\toprule\noalign{}
Stipulated truth & Criterion & Power \\
\midrule\noalign{}
\endfirsthead
\endhead
\bottomrule\noalign{}
\endlastfoot
0.474 & $\hat{\delta} \geq 0.474$ (point-estimate) & \textbf{0.491} \\
0.474 & 95\% CI lower \textgreater{} 0 (any-effect) & 0.868 \\
\end{longtable}
}

Even at $\delta_{\mathrm{truth}} = 0.474$, the (12, 48) design returns $\hat{\delta} \geq
0.474$ in only \textasciitilde49\% of replications, this
\emph{supports} the §4.3 framing that ``not met'' is point-estimate, not
effect-size.

\subsection{D.4 Friedman per-class, discussion of small-N concordance}\label{d.4-friedman-per-class-discussion-of-small-n-concordance}

Per-class Friedman (3 PUTs $\times$ 5 MPs) with Bonferroni $\times$ 4: a chi\^{}2=4.00
raw 0.406 adj 1.000, W=0.333; b chi\^{}2=10.78 raw \textbf{0.029} adj
\textbf{0.116}, W=\textbf{0.898}; c chi\^{}2=4.00 raw 0.406 adj 1.000,
W=0.333; d chi\^{}2=5.00 raw 0.287 adj 1.000, W=0.417. After correction,
no per-class result remains significant at family-wise alpha = 0.05.
Even b-class W = 0.898 (Cohen 1988 large concordance) is insufficient at
N = 3 per class. \textbf{H3 verdict rests on §4.5 sign test}; per-class
Friedman is sensitivity / descriptive only. Friedman main effect on full
60 cells (chi\^{}2 = 15.30, p = 0.0041) confirms MP rank differences but
is logically independent from H3 cross-class consistency.

\textbf{Mixed-effects unavailability.} Primary
\texttt{sms\ \textasciitilde{}\ C(class)\ +\ C(operator)\ +\ C(class):C(operator)\ +\ (1\ \textbar{}\ put)}
returned Singular (insufficient column rank for class $\times$ operator
interaction; N=60 too small for 11-d fixed-effects + 12 PUT random
intercepts). Fallback
\texttt{sms\ \textasciitilde{}\ C(class)\ +\ C(operator)\ +\ (1\ \textbar{}\ put)}
had PUT random-intercept variance hit boundary 0 (degenerates to OLS).
Fallback class p-values (descriptive only): b vs a 0.275 / c vs a 0.892
/ d vs a 0.991. RQ5 conclusion shifts to four-piece presentation (class
means + sign test + Friedman + forest plot).

\subsection{D.5 Pattern-coverage operationalisation}\label{d.5-pattern-coverage-operationalisation}

Per PUT, (MP\_k, R\_outcome $\in$ \{True, False\}) binary-tuple coverage: 5
MPs $\times$ 2 outcomes = 10 cells, PC = \#triggered / 10. Simplest baseline
(per-PUT granularity, not distinguishing mutants). PC range over 12
PUTs: {[}0.500, 1.000{]}, mean 0.733 (\texttt{paper\_numbers\_v4.json}).
Pairing per PUT: Spearman rho = 0.163 (p = 0.613); Kendall tau = 0.136
(p = 0.568).

\textbf{Power caveat.} At n = 12, Spearman 95\% CI $\approx$ {[}-0.5, +0.6{]}
(Fisher-z approximation); cannot distinguish ``zero correlation'' from ``moderate positive'' or
``moderate negative''. p = 0.61 / 0.57 means ``no correlation detected
at n = 12'', not ``correlation does not exist''.

\textbf{Within-class descriptive.} b2 PC = 1.0 with mean SMS = 0.067; b1
PC = 0.7 with mean SMS = 0.20. c3 PC = 1.0 with mean SMS = 0.14; c1 / c2
PC = 0.7-0.8 with mean SMS = 0.0. Within the same class, higher-PC PUT
has lower SMS, contradicts naive ``more PC kills more mutants''
assumption. This pattern is consistent with orthogonality, but n = 12
cannot confirm it. A confirmatory study would need n $\geq$ 30 and a PC
definition that incorporates the mutant dimension.

% [thematic break removed]

\section{E. Stakeholder Deployment Considerations}\label{e.-stakeholder-deployment-considerations-single-output-kernels-scope}

\subsection{E.1 Air-gap incompatibility, full justification}\label{e.1-air-gap-incompatibility-full-justification}

The §5.4 workflow depends on external LLM API calls (Claude / GPT /
DeepSeek), incompatible with most regulated air-gapped V\&V workflows:

{%
\begin{longtable}[]{@{}
  >{\raggedright\arraybackslash}p{(\linewidth - 4\tabcolsep) * \real{0.3333}}
  >{\raggedright\arraybackslash}p{(\linewidth - 4\tabcolsep) * \real{0.3333}}
  >{\raggedright\arraybackslash}p{(\linewidth - 4\tabcolsep) * \real{0.3333}}@{}}
\caption{Air-gap incompatibility, full justification.}\label{tab:p2-25}\\
\toprule\noalign{}
\begin{minipage}[b]{\linewidth}\raggedright
Domain
\end{minipage} & \begin{minipage}[b]{\linewidth}\raggedright
Standard
\end{minipage} & \begin{minipage}[b]{\linewidth}\raggedright
Air-gap requirement
\end{minipage} \\
\midrule\noalign{}
\endfirsthead
\endhead
\bottomrule\noalign{}
\endlastfoot
Nuclear & IEC 60880 & Air-gapped build for safety-critical software \\
Aerospace & DO-178C & Tool-qualified offline build for DAL-A/B; LLM API
outside qualified-tool envelope \\
Medical & IEC 62304 & Class C software lifecycle requires offline
reproducible build \\
Automotive & ISO 26262 & ASIL-D requires offline tool chain; cloud LLM
non-compliant \\
\end{longtable}
}

\textbf{Mitigation paths.} (a) Self-hosted open-weight LLMs (Llama
/ DeepSeek local inference), capability gap vs API-served frontier
models; quantitative impact on SMS untested. (b) Offline-cached
pre-generated mutant pools with commit-hash + raw-response signature
locking, adopters reproduce a published pool offline (this paper's
package supports this), but generating new pools requires external LLM
access. SMS for air-gapped industrial deployment is a future research
direction.

\subsection{E.2 Quarterly batch audit workflow}\label{e.2-quarterly-batch-audit-workflow}

The original per-PR GitHub Actions YAML was removed because its PR-CI
threshold 0.10 contradicted the Appendix~E.3 audit threshold 0.20 and propagated
stakeholder-facing CI assumptions into adopters' pipelines.

\textbf{Replacement: quarterly batch audit:}

{%
\begin{longtable}[]{@{}
  >{\raggedright\arraybackslash}p{(\linewidth - 4\tabcolsep) * \real{0.3333}}
  >{\raggedright\arraybackslash}p{(\linewidth - 4\tabcolsep) * \real{0.3333}}
  >{\raggedright\arraybackslash}p{(\linewidth - 4\tabcolsep) * \real{0.3333}}@{}}
\caption{Quarterly batch audit workflow.}\label{tab:p2-26}\\
\toprule\noalign{}
\begin{minipage}[b]{\linewidth}\raggedright
Step
\end{minipage} & \begin{minipage}[b]{\linewidth}\raggedright
Description
\end{minipage} & \begin{minipage}[b]{\linewidth}\raggedright
Cost
\end{minipage} \\
\midrule\noalign{}
\endfirsthead
\endhead
\bottomrule\noalign{}
\endlastfoot
1 & Offline mutant-pool generation (per PUT 24-30 $\times$ 12 PUTs $\approx$ 300) & LLM
API \$5-15; \textasciitilde30 min \\
2 & \texttt{sms\_campaign.py\ -\/-track\ 2} & 4-core laptop
\textasciitilde10-20 min \\
3 & \texttt{run\_lrca.py} & \textless{} 1 min \\
4 & MR-design team review of MRs with aligned-cell SMS significantly
below historical median (0.275) & 1-2 h meeting \\
\end{longtable}
}

Quarterly: API \textasciitilde\$10 + automation 30 min + manual 1-2 h $\approx$
0.5 person-day. Affordable. Per-PR gating not recommended (LLM latency
5-30 s + cost + non-determinism + air-gap incompatibility); adopters
needing per-PR should use coverage / unit-test gates. SMS does not
replace domain-expert MR physical-reasonableness judgement nor
product-requirements review.

\subsection{E.3 V\&V documentation, conceptual complementarity with ASME V\&V 20-2009}\label{e.3-vv-documentation-conceptual-complementarity-with-asme-vv-20-2009}

This appendix subsection makes \textbf{no normative claim} toward NRC,
FDA, or ISO 26262 review teams. Within single-output kernels scope, no
traceable mapping exists to current IEC 60880 / ISO 26262 / DO-178C /
ASME V\&V 20-2009 normative bodies.

V\&V documentation for scientific computing software (per ASME V\&V
20-2009 §3 and similar guides) requires quantifiable test-adequacy
evidence; current documentation often relies on code coverage + MR lists
+ SME signatures. SMS may appear as \textbf{research-grade supplementary
evidence} alongside coverage and MR lists: (a) aligned-cell SMS per
critical PUT; (b) LRCA three-layer diagnosis (C1 readout / C2-C5
attribution); (c) 60-cell matrix visualisation. Reviewers can
independently run \texttt{REPRODUCIBILITY.md}.

Original threshold recommendations (aligned-cell SMS $\geq$ 0.20 / 0.30 +
C1\_share $\leq$ 0.20) were removed: no normative backing in IEC / ISO /
ASME; misreads as enforce-ready. \textbf{SMS reported as descriptive
supplementary evidence}, with adequacy judgements left to V\&V reviewers
case-by-case.

Conceptually complementary to ASME V\&V 20-2009 §3 code verification
(numerical-solver correctness) vs SMS (MR-set fault-detection adequacy).
Substantial scale gap to multi-module CFD codes; incorporation into V\&V
standards body would require large-scale empirics + multi-year
dialogue with ASME V\&V 20 / IEEE 1012 / IEC 60880 committees.
\textbf{No advocacy that SMS enter any normative certification system in
2027.} SMS does not replace SME signatures, FMECA, PIRT, system-level
V\&V, or ASME V\&V 20 §3 numerical-solver verification, it is one
link in the evidence chain, not a single-point qualification.

% [thematic break removed]

\section{F. Threats to Validity, Detailed Mitigation}\label{f.-threats-to-validity-detailed-mitigation}

\subsection{F.1 Internal threats}\label{f.1-internal-threats}

{%
\begin{longtable}[]{@{}
  >{\raggedright\arraybackslash}p{(\linewidth - 6\tabcolsep) * \real{0.2500}}
  >{\raggedright\arraybackslash}p{(\linewidth - 6\tabcolsep) * \real{0.2500}}
  >{\raggedright\arraybackslash}p{(\linewidth - 6\tabcolsep) * \real{0.2500}}
  >{\raggedright\arraybackslash}p{(\linewidth - 6\tabcolsep) * \real{0.2500}}@{}}
\caption{Internal threats.}\label{tab:p2-27}\\
\toprule\noalign{}
\begin{minipage}[b]{\linewidth}\raggedright
ID
\end{minipage} & \begin{minipage}[b]{\linewidth}\raggedright
Threat
\end{minipage} & \begin{minipage}[b]{\linewidth}\raggedright
Mitigation
\end{minipage} & \begin{minipage}[b]{\linewidth}\raggedright
Residual
\end{minipage} \\
\midrule\noalign{}
\endfirsthead
\endhead
\bottomrule\noalign{}
\endlastfoot
R1 & LLM generation reproducibility & §C.2 reproducibility parameters in
package; dual-LLM cross-source + 20\% manual; reproducibility rests on
archived raw responses and frozen mutant pools (no user-facing seed) & LLM training data may be updated
post-submission \\
R2 & Probabilistic equiv (K\_eq=1000) & §2.3 declares probabilistic
approximation; Hoeffding-style false-equiv bound & K\_eq sweep $\in$
\{500,1000,2000\} deferred to future work \\
R3 & Shared-infrastructure dependency & AVP version pinned to the upstream commit hash;
package embeds AVP source; Li et al. (arXiv) is the infrastructure technical report &
None within scope \\
R4 & LRCA multi-label boundary & §A.2 priority
C5 \textgreater{} C4 \textgreater{} C3 \textgreater{} C2 \textgreater{} C1;
multi-label co-occurrence table in package; root\_cause is ``likely root
cause'' not definitive & C2-C5 may multi-trigger on B/D classes \\
R9 & Mutant-pool size & Pool expanded to mean 17.4: delta moved 0.321 $\to$
0.323, CI narrowed; \textbf{effect size is intrinsic ceiling, not pool
dilution} & Some PUTs \textless{} 30 (cache-bound) \\
R10 & LLM non-determinism & Multi-turn de-dup; K=10/20 repetitions; raw
prompt+response committed for direct reuse & Claude Opus subscription
lacks seed control \\
\textbf{R13} & \textbf{Protocol-implementation gap v3 vs v4} & v4 used 3
providers (Claude 101 / GPT 98 / DeepSeek 99) without dual-blind; v4
C1\_share 0.209 \textgreater{} v3 0.164 weakly opposes ``v4 quality
decline'' & Complete separation requires a follow-up dual-blind rerun on
the v4 grid \\
\end{longtable}
}

R13 is the protocol-asymmetry threat (dual-blind vs V1-V4 only) on the
v3 $\to$ v4 source-axis contrast ($\Delta\delta = -0.009$).

\subsection{F.2 External / construct / conclusion threats}\label{f.2-external-construct-conclusion-threats}

{%
\begin{longtable}[]{@{}
  >{\raggedright\arraybackslash}p{(\linewidth - 6\tabcolsep) * \real{0.2500}}
  >{\raggedright\arraybackslash}p{(\linewidth - 6\tabcolsep) * \real{0.2500}}
  >{\raggedright\arraybackslash}p{(\linewidth - 6\tabcolsep) * \real{0.2500}}
  >{\raggedright\arraybackslash}p{(\linewidth - 6\tabcolsep) * \real{0.2500}}@{}}
\caption{External / construct / conclusion threats.}\label{tab:p2-28}\\
\toprule\noalign{}
\begin{minipage}[b]{\linewidth}\raggedright
ID
\end{minipage} & \begin{minipage}[b]{\linewidth}\raggedright
Threat
\end{minipage} & \begin{minipage}[b]{\linewidth}\raggedright
Class
\end{minipage} & \begin{minipage}[b]{\linewidth}\raggedright
Mitigation
\end{minipage} \\
\midrule\noalign{}
\endfirsthead
\endhead
\bottomrule\noalign{}
\endlastfoot
R5 & 12-PUT representativeness & External & follows the shared site selection;
open class principle (cls extensible); scaling to MD/QC/CFD \\
R6 & Cross-class statistical power & Conclusion & H3 framed exploratory;
mixed-effects unavailable (Singular at N=60); shifted to class means +
sign test + Friedman + forest plot \\
R7 & LLM homogeneity bias & Construct & §C.2 three-LLM rotation +
PUT-class subdivision + 20\% manual sampling; third-party LLM family
validation deferred \\
R8 & LLM-source distributional shift & External & Hit distribution
DeepSeek 11/15, Claude 4/15, GPT 0/15; §3.5 argument is categorical
AST-locality, not hit frequency; HP/SI/TF zero-overlap across all 3
sources unaffected \\
R12 & HOM equivalence & External & Tool-unreachability claim confined to
first-order syntactic tools; HOM empirical testing listed as future
work \\
\end{longtable}
}

Conclusion threats: multiple comparisons handled by BH-FDR at
$\alpha_{\mathrm{FDR}}=0.05$ (§D.1); N=20 stability handled by 1,000-iteration
bootstrap CI (§D.3). Construct threat, SMS measures epistemological
semantic detection, not production engineering value (out of scope here).

% [thematic break removed]

\section{G. SMS-to-MS Degeneration Theorem, Full Proof}\label{g.-sms-ms-degeneration-theorem-full-proof}

\subsection{G.1 Notation cross-reference}\label{g.1-notation-cross-reference-with-2.1}

PUT $S_i$, mutation operator family mut (syntactic or semantic),
mutant set $\mathrm{mut}(S_i)$, MP set, equivalence tolerance $\varepsilon_{\mathrm{eq}}$, equivalence
sample size $K_{\mathrm{eq}}$, AVP tolerance $\varepsilon_{\mathrm{AVP}}$. Three-state decomposition:
$\mathrm{mut}(S)=\mathrm{killed} \cup \mathrm{equiv} \cup \mathrm{survive}$ (disjoint). SMS
formula:

\[
\mathrm{SMS}_{i,k,j} = \frac{|\mathrm{killed}_{i,k,j}|}{|\mathrm{mut}_j(S_i)| - |\mathrm{equiv}_{i,k,j}|}
\]

\subsection{G.2 Degenerate-limit definition}\label{g.2-degenerate-limit-definition-r-8-p1-3-revision-3-joint-conditions}
% NOTE: D_S -> \mathcal{D}_P rename is scoped to the degeneration-theorem sections in this pass; the global sweep completes in the notation-alignment pass (T6.1).

The degenerate limit $L = L_{\mathrm{equiv}} \wedge L_{\mathrm{killed}} \wedge L_{\mathrm{mut}}$
consists of three \textbf{joint conditions} (paired axes), each
controlling one layer of the SMS formula. L1-L6 are not 6 independent
axes; the pairing reflects dependency structure across the SMS formula
layers. Equivalently, the same conjunction regroups orthogonally as
$L = L_{\mathrm{lim}} \wedge L_{\mathrm{switch}}$, where
$L_{\mathrm{lim}} = L1 \wedge L2 \wedge L3$ collects the quantitative
limit axes ($\varepsilon_{\mathrm{eq}} \to 0$,
$K_{\mathrm{eq}} \to \infty$, $\varepsilon_{\mathrm{AVP}} \to 0$) and
$L_{\mathrm{switch}} = L4 \wedge L5 \wedge L6$ collects the discrete
configuration switches (MP set $= \{\mathrm{MP}_{\mathrm{eq}}\}$;
rule-based syntactic operators; imperative deterministic PUT class).

{%
\begin{longtable}[]{@{}
  >{\raggedright\arraybackslash}p{(\linewidth - 4\tabcolsep) * \real{0.3333}}
  >{\raggedright\arraybackslash}p{(\linewidth - 4\tabcolsep) * \real{0.3333}}
  >{\raggedright\arraybackslash}p{(\linewidth - 4\tabcolsep) * \real{0.3333}}@{}}
\caption{Degenerate-limit definition.}\label{tab:p2-29}\\
\toprule\noalign{}
\begin{minipage}[b]{\linewidth}\raggedright
Joint condition
\end{minipage} & \begin{minipage}[b]{\linewidth}\raggedright
Component axes
\end{minipage} & \begin{minipage}[b]{\linewidth}\raggedright
Pairing rationale
\end{minipage} \\
\midrule\noalign{}
\endfirsthead
\endhead
\bottomrule\noalign{}
\endlastfoot
\textbf{$L_{\mathrm{equiv}}$} (Lemma 9.1) & L1: $\varepsilon_{\mathrm{eq}} \to 0$; L2: $K_{\mathrm{eq}} \to \infty$ & L1 alone:
equiv stays probabilistic ($K_{\mathrm{eq}}$ cannot cover $\mathcal{D}_P$). L2 alone: bitwise
equality diluted by $\varepsilon_{\mathrm{eq}}$. Both simultaneously degenerate equiv to
classical behavioural equivalence outside a $\mathcal{D}_P$-measure-zero set. \\
\textbf{$L_{\mathrm{killed}}$} (Lemma 9.2) & L3: $\varepsilon_{\mathrm{AVP}}^k \to 0$; L4: MP =
$\{\mathrm{MP}_{\mathrm{eq}}\}$, R(y,y') $\equiv$ y = y' & L3 alone: $\varepsilon_{\mathrm{AVP}} \to 0$ still allows
non-trivial MP (monotonicity, convergence). L4 alone: equality still
carries $\varepsilon_{\mathrm{AVP}}$ tolerance. Both simultaneously degenerate killed to
classical difference detection. \\
\textbf{$L_{\mathrm{mut}}$} (Lemma 9.3) & L5: $\mathrm{mut}_j \to$ rule-based syntactic (Mothra
AOR/ROR/SDL/CRP); L6: cls(I) $\subseteq$ \{imperative deterministic\} & L5 alone:
syntactic ops on probabilistic/ML may still trigger domain-semantic
subsets. L6 alone: imperative programs still mutable by OS/HP/TF/SI.
Both simultaneously degenerate mut(S) to Jia \& Harman syntactic mutant
set. \\
\end{longtable}
}

\subsection{G.3 Lemmas, Three-state decomposition degenerates under L}\label{g.3-lemmas-three-state-decomposition-degenerates-under-l}

\textbf{Lemma 9.1 (equiv degeneration; exception-set
qualification).} Assume
$\mathrm{supp}(\mathcal{D}_P) \supseteq \mathcal{X}_{\mathrm{adm}}$.
Under $L_{\mathrm{equiv}}$ (L1 $\wedge$ L2), semantic-class equivalence
(E1 $\wedge$ E2) degenerates to classical behavioural equivalence
\textbf{outside an exception set $N$}: over floating-point (finite)
input domains, $N$ is the finite set of pathological inputs
(Not-a-Number propagation and overflow/underflow traps), empty in the
exact discrete case; under the continuum idealisation, $N$ has
$\mathcal{D}_P$-measure zero.

\emph{Proof.} E1 (type consistency) holds trivially under $\varepsilon_{\mathrm{eq}} \to 0$ (L6
makes imperative output spaces scalar/vector with static types). E2
($|S_i(x) - s'(x)| < \varepsilon_{\mathrm{eq}}$ for $K_{\mathrm{eq}}$ samples)
under L1 ($\varepsilon_{\mathrm{eq}} \to 0$) $\wedge$ L2 ($K_{\mathrm{eq}} \to \infty$ with measure-equivalent
sampling, which by the support assumption reaches all of
$\mathcal{X}_{\mathrm{adm}}$)
is equivalent to $\forall x \in \mathcal{X}_{\mathrm{adm}} \setminus N: S_i(x) = s'(x)$, where
$N$ is the exception set of the lemma statement (floating-point
domains: the finite set of pathological inputs / NaN propagation;
exact discrete case: $N = \emptyset$; continuum idealisation:
$\mathcal{D}_P$-measure zero). This
matches \citet{d7c38286734419b52de4262c9802ebdfcf4b9447} §3 classical equivalent-mutant definition
up to the exception set $N$.

\textbf{Lemma 9.2 (killed degeneration).} Under L3 $\wedge$ L4, killed
degenerates to classical difference detection.

\emph{Proof.} L4 restricts MP to \{MP\_eq\} with R(y, y') $\equiv$ y = y'. For
mr = (r, R) and mutant s', the $\mathrm{MP}_{\mathrm{eq}}$ violation condition $\exists x: S_i(x) \neq
s'(r(x))$ under L3 ($\varepsilon_{\mathrm{AVP}} \to 0$) becomes exact inequality. With r = id,
violation is $S_i(x) \neq s'(x)$, classical difference detection. With r
$\neq$ id, $\mathrm{MP}_{\mathrm{eq}}$ still requires $S_i(x) = s'(r(x))$ as reference oracle from
original program; no new state classifications introduced.

\textbf{Lemma 9.3 (mut degeneration).} Under L5 $\wedge$ L6, mut\_j(S\_i)
degenerates to the syntactic mutant set of \citet{d7c38286734419b52de4262c9802ebdfcf4b9447}.

\emph{Proof.} L5 switches $\mathrm{mut}_j$ to rule-based syntactic operators (AOR,
ROR, SDL, CRP, UOI, standard Mothra/Proteum sets); L6 restricts PUTs
to imperative deterministic programs, excluding triggering conditions
for semantic operators on probabilistic/ML programs. Under this
configuration, $\mathrm{mut}_j(S_i)$ is the literature-defined syntactic mutant
set, independent of domain semantics.

\subsection{G.4 Theorem 9.1, Detailed proof}\label{g.4-theorem-9.1-detailed-proof}

\textbf{Theorem 9.1 (SMS $\to$ MS degeneration).} Fix the switch
component $L_{\mathrm{switch}} = L4 \wedge L5 \wedge L6$ and assume
$\mathrm{supp}(\mathcal{D}_P) \supseteq \mathcal{X}_{\mathrm{adm}}$.
Along the limit component $L_{\mathrm{lim}} = L1 \wedge L2 \wedge L3$,
outside the exception set $N$ of Lemma 9.1 (finite set of pathological
inputs over floating-point domains, empty in the exact discrete case;
$\mathcal{D}_P$-measure-zero under the continuum idealisation),

\[ \mathrm{SMS}_{i,k,j} \xrightarrow{L_{\mathrm{lim}}} \mathrm{MS}_{i,j} := \frac{|\mathrm{killed}_{i,j}^{\mathrm{classic}}|}{|\mathrm{mut}_j^{\mathrm{syntax}}(S_i)| - |\mathrm{equiv}_{i,j}^{\mathrm{classic}}|} \]

\emph{Proof.} Combine Lemmas 9.1-9.3:

\begin{itemize}

\item
  Numerator: under L3 $\wedge$ L4, $\mathrm{killed}_{i,k,j} \to
  \mathrm{killed}_{i,j}^{\mathrm{classic}}$ (Lemma 9.2); L4 makes $\mathrm{MR}_{i,k}$ trivial
  over k (only $\mathrm{MP}_{\mathrm{eq}}$ remains), so subscript k degenerates.
\item
  Denominator $|\mathrm{mut}_j(S_i)|$: under L5 $\wedge$ L6 $\to$
  $|\mathrm{mut}_j^{\mathrm{syntax}}(S_i)|$ (Lemma 9.3).
\item
  Denominator $|\mathrm{equiv}_{i,k,j}|$: under L1 $\wedge$ L2 $\to$
  $|\mathrm{equiv}_{i,j}^{\mathrm{classic}}|$ (Lemma 9.1;
  outside the exception set $N$).
\end{itemize}

Substituting into the SMS formula gives $\mathrm{MS}_{i,j}$.

\textbf{Corollary 9.1 (LRCA trivialisation).} Under $L$, C
= \{C1, \ldots, C5\} degenerates to \{C1\}.

\emph{Sketch.} Each of C2-C5's triggering precondition depends on at
least one $L_j$ being violated. When the three joint conditions
$L_{\mathrm{equiv}} \wedge L_{\mathrm{killed}} \wedge L_{\mathrm{mut}}$ hold simultaneously,
every non-trivial-space dimension (MP non-triviality, AVP tolerance
non-zero, non-empty equiv set, MR-design DOF, class-mapping openness)
closes; C2-C5 triggering set becomes empty (read off from §A.2 decision
tree). The per-C\_k to per-L\_j minimum-sufficient mapping depends on
§3.11 LRCA-classifier engineering thresholds; we do not claim one-to-one
correspondence at formal level, readers can trace through §A.2 + §C.4.
Under L, $\texttt{suspect\_share} \to 0$, LRCA reports only C1, SMS degenerates to
single-layer metric consistent with \citet{d7c38286734419b52de4262c9802ebdfcf4b9447} MS engineering
structure.

\textbf{Empirical consistency.} Theorem 9.1 + Corollary 9.1 jointly
guarantee any SMS-based empirical conclusion (Cliff's delta §4.3,
Friedman chi\^{}2 §4.5, Spearman rho §4.6) is structurally consistent
with existing \citet{d7c38286734419b52de4262c9802ebdfcf4b9447} literature in classical
syntactic-mutation scenarios, no metric-level semantic fragmentation.

\section{H. Adjoint Extension Arm, Subjects, Pools, and Forward
Results}\label{h.-adjoint-extension-arm}

This appendix gives the full specification of the confirmatory adjoint
extension arm summarised in §3.7 and §4.9 of the main paper. The arm
exercises the adjoint invariant $\psi_6$ on three third-party
linear-operator libraries, disjoint from the 12 PUTs and the
pre-registered 60-cell. Each library carries a real fixed defect in its
adjoint code path as a ground-truth anchor:

\begin{itemize}

\item
  \textbf{scipy.sparse.linalg LinearOperator} (linear algebra; BSD-3).
  Program under test: the scaled-operator adjoint $M^{\mathrm{H}}$ for
  $M = \alpha A$ with complex scale $\alpha$. Real anchor: scipy issue
  \#8900 / PR \#8962 (the adjoint dropped the conjugate on $\alpha$,
  returning $\alpha A^{\mathrm{H}}$ instead of $\bar{\alpha}
  A^{\mathrm{H}}$).
\item
  \textbf{pylops.signalprocessing.FFT} with \texttt{real=True} (signal
  processing; BSD-3). Program under test: the real-FFT adjoint
  \texttt{\_rmatvec}; the adjoint identity is the library's own
  \emph{dot-test}. Real anchor: pylops issue \#268 / PR \#292 (commit
  \texttt{bd3a919}), a wrong scaling of the middle/Nyquist frequency bin
  that made the adjoint inconsistent with the forward.
\item
  \textbf{jax} \texttt{linear\_transpose} / \texttt{vjp} (automatic
  differentiation; Apache-2.0). Program under test: the autodiff-derived
  transpose $M^{\mathrm{T}}$ of a linear operator containing an inverse
  real FFT, $M(x) = \mathrm{irfft}(\mathrm{rfft}(x)\cdot w)$. Real anchor:
  jax issue \#6223 / PR \#6232 (commit \texttt{112ae41}), an
  \texttt{irfft} transpose rule that reversed order with an off-by-one
  index, so \texttt{vjp} silently produced a non-adjoint transpose.
\end{itemize}

For each subject the MR families are adj.self (self-adjointness) and
adj.dual (scaled/dual-operator transpose). The mutant pool spans
adjoint-breaking fault mechanisms (drop-conjugate, sign flip, scale and
phase errors, transpose-without-conjugation, Nyquist mis-scaling), with
the real fixed defect included among them, plus semantically equivalent
mutants (scalar-multiplication reorder, conjugation of a real adjoint
output) that a strong MR must \emph{not} kill. Because the pre-fix
releases are uninstallable on the host platform (no \texttt{arm64} wheel),
each real defect is reproduced in extracted-diff form on the current
release, mirroring the boundary study; $\varepsilon_{\mathrm{tol}}$ is an
explicit factor throughout.

{%
\begin{longtable}[]{@{}p{2.0cm}p{3.2cm}p{2.5cm}p{1.8cm}p{1.8cm}@{}}
\caption{Adjoint extension arm, forward kill.}\label{tab:suppl-adjoint}\\
\toprule\noalign{}
Substrate & Operator class (real anchor) & Active killed / total & Equiv.\ survived & Forward SMS \\
\midrule\noalign{}
\endfirsthead
\endhead
\bottomrule\noalign{}
\endlastfoot
scipy LinearOperator & linear algebra (\#8900) & 5 / 5 & 2 / 2 & \textbf{1.00} \\
pylops FFT real & signal processing (\#292) & 4 / 4 & 2 / 2 & \textbf{1.00} \\
jax vjp / transpose & automatic differentiation (\#6223) & 4 / 4 & 2 / 2 & \textbf{1.00} \\
\end{longtable}
}

Across all three substrates the correct program survives, every
non-equivalent (active) mutant is killed, including the three real
fixed defects, and no semantically equivalent mutant is killed, so the
forward SMS is 1.00 with zero false positives on the equivalent set. The
scipy \#8900 defect is shared with the §4.8 boundary study (adj.self case)
and is not counted as independent corroboration twice; the pylops \#292
and jax \#6223 defects are unique to this arm. The pools are
hand-constructed and low-cardinality, so the arm demonstrates the duality
forward rather than measuring a high-adequacy MR set on a discovered fault
population (main-paper threat R14).

\section{I. Result-Level Real-Defect Evidence Ledger}\label{i.-result-level-real-defect-evidence-ledger}

This appendix records only the result-level facts used by the main
manuscript's real-defect evidence slice. It deliberately excludes
candidate mining, rejected cases, repository tooling, and benchmark-release
narrative. The purpose is auditability for the semantic-mutation argument:
kill-rate, semantic-stratum alignment, and real-defect detection are related
but distinct constructs.

{%
\small
\begin{longtable}[]{@{}p{0.18\linewidth}p{0.52\linewidth}p{0.24\linewidth}@{}}
\caption{Result-level real-defect evidence ledger.}\label{tab:realdefect-ledger}\\
\toprule\noalign{}
Audit item & Ledger content used in the manuscript & Main-text role \\
\midrule\noalign{}
\endfirsthead
\endhead
\bottomrule\noalign{}
\endlastfoot
Verified stratum counts & 34 verified MR-detectable defects across five
meta-patterns: \(m_{\mathrm{mono}}\) 10, \(m_{\mathrm{conv}}\) 9,
\(m_{\mathrm{inv}}\) 6, \(m_{\mathrm{adj}}\) 5, and
\(m_{\mathrm{rev}}\) 4 & Shows that the declared semantic strata are not
only toy labels \\
Mutation-phase aggregate & T1--B1 mean difference +0.101,
Holm-adjusted \(p = 0.046\), Cliff's \(\delta = +0.247\); T1--A1 not
significant; B1--B2 not significant & Supports a modest, qualified
mutation-score advantage rather than a dominance claim \\
Real-defect face & T1 detects 30/30 tabled real defects; B1 has nonzero
detection in 4/30, B2 in 6/30; A1-a loses 17/30, A1-b loses 16/30, and
both lose 10/30 & Separates real-defect detection from aggregate mutant
kill-rate \\
Non-nesting cases & A-LAPACK-004, A-OPENBLAS-001, B-POCKETFFT-002, and
E-ORDINARYDIFFEQ-001 & Supports the semantic-effect fiber view rather
than a single scalar hierarchy \\
\end{longtable}
\normalsize
}

The ledger is intentionally a result table, not a corpus or benchmark
description. The main manuscript uses it only to support construct
separation: a relation family can score similarly on aggregate mutant
kill-rate while differing sharply on real-defect detection, and different
relation families can observe non-nested semantic-effect fibers.

\bibliographystyle{ACM-Reference-Format}
\bibliography{references}

\end{document}
